As árvores CTT refinam a HTA congelada \cite[pp.~3, 5, 8--9]{hts} por responsabilidade e ordem temporal \cite[cap.~4, pp.~39--66]{paterno}, sem alterar objetivos, operações nem critérios de sucesso. A identifica tarefa abstrata, U tarefa do usuário (operação não observável), I tarefa interativa (gera evento de interação) e S tarefa do sistema (gera evento do sistema). $\gg$ impõe sequência; $[\,]\gg$ transmite o resultado ao sucessor; $[\,]$ oferece escolha; $|||$ permite intercalação; $|=|$ permite qualquer ordem sem intercalação. Colchetes em torno de uma tarefa indicam opcionalidade e o asterisco indica repetição. Os números recuperam as operações da HTA.
\par Para cada raiz $T$ aplica-se $(T\,|>\,\text{consultar ajuda})\,[>\,(\text{desistir}\,[\,]\,\text{atingir limite}\,[\,]\,\text{falhar})$. A consulta suspende e retoma: solicitar ajuda (I) $\gg$ apresentar dica (S) $\gg$ [escolher camada (I) $\gg$ apresentar conteúdo (S)]* $\gg$ fechar (I) $\gg$ restaurar origem (S). A desativação termina por desistência (I), limite (S) ou falha técnica (S). A leitura mental não produz evento; a exposição do signo registra uma oportunidade de leitura. Em C2, o início mensurado tem o pedido carregado e um acesso à revisão, o que conserva as três ações de referência (HO). V1 conserva as três categorias permutáveis, com etiquetas $|=|$ preços, e a repetição em caso de divergência. V2 conserva a espera pelo carregamento, a reconsulta opcional e a remoção do residual. Prevenção e tratamento de erros ficam na matriz de rupturas, fora do CTT.
